% !TeX program = xelatex
% =====================================================================
%  main.tex
%  Template de apresentacao para o Bacharelado em Fisica
%  Beamer + XeLaTeX/LuaLaTeX + modulos Lua (pasta ./lua)
%
%  Compile com:  latexmk -xelatex main.tex   (funciona no UaiTeX)
%                latexmk -lualatex main.tex (usa os modulos Lua)
%  ou:           make
% =====================================================================
\RequirePackage{iftex}
\ifPDFTeX
  \errmessage{%
    Este template exige XeLaTeX ou LuaLaTeX (fontspec, unicode-math).
    Use: latexmk -xelatex main.tex  ou  make}%
  \expandafter\endinput
\fi

\documentclass[aspectratio=169,10pt]{beamer}

\usepackage{babel}
\babelprovide[import,main]{brazilian}
\usetheme{fisica}
\usepackage{fisica}
\usepackage{fisica-video}

% ---------------------------------------------------------------------
%  Metadados (edite esta parte)
% ---------------------------------------------------------------------
\title[Oscilador quântico]{O oscilador harmônico quântico}
\subtitle{Da energia de ponto zero aos qubits}

\author[Seu Nome]{\texorpdfstring{Seu Nome\\[2pt]
  {\small Orientador: Prof.\ Dr.\ Fulano de Tal}}{Seu Nome}}
\institute[UFXX]{Universidade Federal de Exemplo\\
  Bacharelado em Física}
\date{\today}

\fisicaeyebrow{FÍSICA \textperiodcentered{} BACHARELADO \textperiodcentered{} SEMINÁRIO}

% =====================================================================
\begin{document}

% ---------------------------------------------------------------------
%  Abertura
% ---------------------------------------------------------------------
\frame[plain]{\titlepage}

\begin{frame}{Roteiro}
  \begin{enumerate}
    \roteiro
  \end{enumerate}
\end{frame}

% Paginas de secao automaticas
\AtBeginSection[]{%
  \begin{frame}[plain]
    \sectionpage
  \end{frame}%
}

% ---------------------------------------------------------------------
%  Conteudo
% ---------------------------------------------------------------------
\section{Motivação}
\begin{frame}{Por que o oscilador harmônico?}
  \begin{columns}[T,onlytextwidth]
    \begin{column}{0.54\textwidth}
      \begin{itemize}
        \item Perto de um mínimo local com \alert{$V''(x_0)>0$},
              o potencial é aproximadamente parabólico:
      \end{itemize}
      \begin{equacao}
        \[
          V(x) \approx V(x_0)
          + \tfrac12 V''(x_0)\,(x-x_0)^2
        \]
      \end{equacao}
      \begin{itemize}
        \item Por isso ele descreve \alert{moléculas}, \alert{sólidos},
              \alert{campos} e \alert{circuitos}.
        \item É um dos poucos modelos quânticos com solução
              \alert{analítica exata}.
      \end{itemize}
    \end{column}
    \begin{column}{0.42\textwidth}
      \begin{destaque}{Onde ele aparece}
        \begin{itemize}
          \item Vibrações moleculares (espectroscopia IV)
          \item Fônons em cristais
          \item Íons aprisionados e qubits
          \item Ressonadores nanomecânicos
        \end{itemize}
      \end{destaque}
    \end{column}
  \end{columns}
\end{frame}

\begin{frame}{Clássico versus quântico}
  \centering
  \footnotesize
  \renewcommand{\arraystretch}{1.35}
  \begin{tabular}{@{}>{\raggedright\arraybackslash}p{0.16\linewidth}
                    >{\raggedright\arraybackslash}p{0.37\linewidth}
                    >{\raggedright\arraybackslash}p{0.41\linewidth}@{}}
    \toprule
    \textbf{Aspecto} & \textbf{Clássico} & \textbf{Quântico} \\
    \midrule
    \comparacao
    \bottomrule
  \end{tabular}

  \vspace{0.7em}
  \begin{destaque}[fisicaamber]{A diferença central}
    No regime quântico a energia não pode ser zero:
    $E_0=\tfrac12\hbar\omega$.
  \end{destaque}
\end{frame}


\section{Formulação quântica}
\begin{frame}{A equação de Schrödinger}
  \begin{equacao}
    \[
      -\frac{\hbar^2}{2m}\,\frac{\mathrm{d}^2\psi}{\mathrm{d}x^2}
      + \frac12\,m\omega^2 x^2\,\psi
      = E\,\psi
    \]
  \end{equacao}

  \begin{columns}[T,onlytextwidth]
    \begin{column}{0.48\textwidth}
      \begin{destaque}{Adimensionalizar}
        \[
          \xi = \sqrt{\frac{m\omega}{\hbar}}\,x,
          \qquad
          \epsilon = \frac{2E}{\hbar\omega}
        \]
        \[
          \frac{\mathrm{d}^2\psi}{\mathrm{d}\xi^2}
          + \left(\epsilon-\xi^2\right)\psi = 0
        \]
      \end{destaque}
    \end{column}
    \begin{column}{0.48\textwidth}
      \begin{destaque}[fisicaviolet]{Condições de contorno}
        \begin{itemize}
          \item $\psi(\xi)\to 0$ quando $|\xi|\to\infty$
          \item $\psi$ e $\psi'$ contínuas
          \item Soluções normalizáveis exigem
                $\epsilon = 2n+1$
        \end{itemize}
      \end{destaque}
    \end{column}
  \end{columns}
\end{frame}

\begin{frame}{Espectro de energia}
  \begin{equacao}
    \[
      E_n = \left(n+\frac12\right)\hbar\omega,
      \qquad n = 0,1,2,\ldots
    \]
  \end{equacao}

  \begin{columns}[T,onlytextwidth]
    \begin{column}{0.53\textwidth}
      \begin{itemize}
        \item Níveis \alert{igualmente espaçados}:
              $\Delta E = \hbar\omega$.
        \item O estado fundamental tem energia \alert{não nula}.
      \end{itemize}
      \begin{destaque}{Energia de ponto zero (via Lua)}
        Para $\hbar\omega = 250\ \mathrm{meV}$:
        \[
          E_0 = \frac12\hbar\omega = \energia{0}{250}
        \]
      \end{destaque}
    \end{column}
    \begin{column}{0.44\textwidth}
      \begin{destaque}[fisicaamber]{Consequência}
        Mesmo a $T=0$ o oscilador não está em repouso. Essa energia
        residual é exigida pelo princípio da incerteza e tem efeitos
        mensuráveis em vibrações moleculares e em qubits.
      \end{destaque}
    \end{column}
  \end{columns}
\end{frame}


\section{Espectro e funções de onda}
\begin{frame}{Funções de onda}
  \framesubtitle{Soluções normalizadas sobre o potencial parabólico}
  \begin{columns}[T,onlytextwidth]
    \begin{column}{0.35\textwidth}
      \begin{destaque}{Forma geral}
        \[
          \psi_n(\xi) = N_n\,H_n(\xi)\,e^{-\xi^2/2}
        \]
        {\scriptsize\color{fisicamuted}%
         $N_n=(2^n n!\sqrt{\pi})^{-1/2}$,\quad
         $\xi=\sqrt{m\omega/\hbar}\,x$}
      \end{destaque}
      \begin{itemize}
        \item $n$ nós e paridade $(-1)^n$
        \item $|\psi_n|^2$: densidade de probabilidade
      \end{itemize}
    \end{column}
    \begin{column}{0.62\textwidth}
      \centering
      \begin{tikzpicture}
        \begin{axis}[
          width=\linewidth,
          height=5.55cm,
          axis lines=left,
          axis line style={draw=fisicamuted,line width=0.5pt},
          xmin=-4.1, xmax=4.1,
          ymin=0, ymax=5.7,
          xlabel={$x\,/\,\sqrt{\hbar/m\omega}$},
          ylabel={$E\,/\,\hbar\omega$},
          xlabel style={font=\scriptsize,text=fisicamuted},
          ylabel style={font=\scriptsize,text=fisicamuted},
          tick label style={font=\scriptsize,text=fisicamuted},
          grid=major,
          grid style={draw=fisicamuted!12},
          legend style={
            at={(0.03,0.97)},anchor=north west,
            draw=none,fill=fisicabg,fill opacity=0.85,text opacity=1,
            text=fisicafg,font=\scriptsize,row sep=0pt},
          legend cell align=left,
        ]
          \plotOscilador
        \end{axis}
      \end{tikzpicture}
    \end{column}
  \end{columns}
\end{frame}

\begin{frame}{Espectro em unidades práticas}
  \begin{columns}[T,onlytextwidth]
    \begin{column}{0.56\textwidth}
      \centering
      \footnotesize
      \renewcommand{\arraystretch}{1.25}
      \begin{tabular}{@{}cccc@{}}
        \toprule
        $n$ & $E_n$ (meV) & paridade & nós \\
        \midrule
        \niveisTabela
        \bottomrule
      \end{tabular}

      \vspace{0.5em}
      {\scriptsize\color{fisicamuted}%
       Gerado por \texttt{lua/conteudo.lua} com
       $\hbar\omega = 250\ \mathrm{meV}$.}
    \end{column}
    \begin{column}{0.41\textwidth}
      \begin{destaque}[fisicagreen]{Leitura da tabela}
        \begin{itemize}
          \item O espaçamento é constante: $250\ \mathrm{meV}$.
          \item A paridade alterna com $n$.
          \item O número de nós cresce linearmente.
        \end{itemize}
      \end{destaque}
    \end{column}
  \end{columns}
\end{frame}

\begin{frame}{Onde isso aparece}
  \begin{columns}[T,onlytextwidth]
    \begin{column}{0.49\textwidth}
      \begin{destaque}{Moléculas diatômicas}
        Vibrações em torno do equilíbrio. As transições
        $\Delta n = \pm 1$ explicam as bandas no infravermelho.
      \end{destaque}
      \begin{destaque}[fisicaamber]{Íons aprisionados}
        O movimento de um íon em um potencial efetivo forma um
        oscilador; seus estados são usados em computação quântica.
      \end{destaque}
    \end{column}
    \begin{column}{0.49\textwidth}
      \begin{destaque}[fisicaviolet]{Circuitos supercondutores}
        Junções Josephson formam um oscilador anarmônico cujos dois
        primeiros níveis definem um qubit.
      \end{destaque}
      \begin{destaque}[fisicagreen]{Nanomecânica}
        Osciladores mecânicos levados ao estado fundamental testam
        os limites da mecânica quântica.
      \end{destaque}
    \end{column}
  \end{columns}
\end{frame}


\section{Aplicações e ferramentas}
\begin{frame}{Constantes físicas via Lua}
  \centering
  \footnotesize
  \renewcommand{\arraystretch}{1.2}
  \begin{tabular}{@{}lll@{}}
    \toprule
    \textbf{Símbolo} & \textbf{Valor (CODATA 2022)} & \textbf{Grandeza} \\
    \midrule
    \tabelaconstantes{c,hbar,e,me,mp,kB}
    \bottomrule
  \end{tabular}

  \vspace{0.8em}
  \begin{destaque}{Como funciona}
    \small
    \texttt{\textbackslash tabelaconstantes\{...\}} lê
    \texttt{lua/constantes.lua} e imprime linhas já formatadas pelo
    \texttt{siunitx}. Alterar um valor no Lua atualiza todos os slides.
  \end{destaque}
\end{frame}

\begin{frame}[fragile]{Como o template usa Lua}
  \begin{columns}[T,onlytextwidth]
    \begin{column}{0.45\textwidth}
      \begin{destaque}{Um módulo, três usos}
        \begin{itemize}
          \item Constantes e unidades
          \item Tabelas e listas de exercícios
          \item CSV $\rightarrow$ gráficos \texttt{pgfplots}
        \end{itemize}
      \end{destaque}
      \begin{destaque}[fisicaviolet]{Compilar}
        \small
        \texttt{latexmk -lualatex main.tex}\\[2pt]
        ou simplesmente \texttt{make}
      \end{destaque}
      {\scriptsize\color{fisicamuted}%
       No texto: \texttt{\textbackslash cte\{hbar\}} e
       \texttt{\textbackslash energia\{0\}\{250\}}.}
    \end{column}
    \begin{column}{0.52\textwidth}
      \begin{lstlisting}[style=lua]
-- lua/fisica.lua (trecho)
function M.constante(chave)
  local c = constantes.obter(chave)
  tex.sprint(string.format(
    "$%s = \\qty{%s}{%s}$",
    c.simbolo,
    c.valor,
    c.unidade))
end
      \end{lstlisting}
    \end{column}
  \end{columns}
\end{frame}

\begin{frame}{Exercícios}
  \small
  \setlength{\itemsep}{0.55em}
  \begin{enumerate}
    \exercicios
  \end{enumerate}

  \vfill
  {\scriptsize\color{fisicamuted}%
   Lista montada a partir de \texttt{lua/conteudo.lua}.}
\end{frame}

\begin{frame}{Resumo}
  \begin{columns}[T,onlytextwidth]
    \begin{column}{0.49\textwidth}
      \begin{destaque}{O que vimos}
        \begin{itemize}
          \item Potencial parabólico $\Rightarrow$ modelo exato
          \item $E_n=\left(n+\tfrac12\right)\hbar\omega$
          \item Funções de onda com paridade definida
        \end{itemize}
      \end{destaque}
    \end{column}
    \begin{column}{0.49\textwidth}
      \begin{destaque}[fisicagreen]{Próximos passos}
        \begin{itemize}
          \item Operadores de criação e destruição
          \item Oscilador em $3$D e degenerescência
          \item Oscilador forçado e dissipação
        \end{itemize}
      \end{destaque}
    \end{column}
  \end{columns}

  \vspace{0.5em}
  \begin{equacao}
    \[
      \hat a = \sqrt{\frac{m\omega}{2\hbar}}
      \left(\hat x + \frac{\iu}{m\omega}\,\hat p\right),
      \qquad
      \comm{\hat a}{\hat a^\dagger} = 1
    \]
  \end{equacao}
\end{frame}

\begin{frame}{Vídeo local e animação}
  \framesubtitle{Arquivos em \texttt{videos/}, sem depender de conexão}
  \begin{columns}[T,onlytextwidth]
    \begin{column}{0.53\textwidth}
      \centering
      \videolocal[width=0.94\linewidth]{demo-oscilador.mp4}{demo-oscilador.png}\\[2pt]
      {\scriptsize\color{fisicamuted}%
       Superposição de $\psi_0$ e $\psi_1$: a densidade $|\psi|^2$
       oscila dentro do potencial.}
    \end{column}
    \begin{column}{0.43\textwidth}
      \begin{destaque}{Comandos}
        {\scriptsize\ttfamily
         \textbackslash videolocal[width=...]\\
         \hspace*{1.1em}\char`\{video.mp4\char`\}\char`\{capa.png\char`\}\\[3pt]
         \textbackslash videolink[texto]\char`\{video.mp4\char`\}}%
      \end{destaque}
      \vspace{0.15em}
      {\small\videolink[Abrir no player do sistema]{demo-oscilador.mp4}}
      \vspace{0.45em}

      {\scriptsize\color{fisicamuted}%
       A reprodução embutida depende do leitor (\texttt{media9});
       o link abre o vídeo no player externo.}
    \end{column}
  \end{columns}
\end{frame}

\begin{frame}{Referências}
  \small
  \begin{thebibliography}{9}
    \bibitem{griffiths}
      D. J. Griffiths,
      \emph{Introduction to Quantum Mechanics},
      2.\ ed., Pearson, 2005.

    \bibitem{sakurai}
      J. J. Sakurai,
      \emph{Modern Quantum Mechanics},
      2.\ ed., Addison-Wesley, 1994.

    \bibitem{cohen}
      C. Cohen-Tannoudji, B. Diu, F. Laloë,
      \emph{Quantum Mechanics}, vol.\ 1, Wiley, 1977.

    \bibitem{codata}
      P. J. Mohr, D. B. Newell, B. N. Taylor,
      \emph{CODATA recommended values of the fundamental physical
      constants: 2022}, NIST.
  \end{thebibliography}
\end{frame}


% ---------------------------------------------------------------------
%  Encerramento
% ---------------------------------------------------------------------
\begin{frame}[plain]
  \centering
  \vfill
  {\usebeamerfont{title}\usebeamercolor[fg]{title}Obrigado!\par}
  \vspace{0.30cm}
  {\color{fisicacyan}\rule{2.2cm}{1.5pt}}\par
  \vspace{0.40cm}
  {\usebeamerfont{subtitle}\usebeamercolor[fg]{subtitle}Perguntas?\par}
  \vfill
\end{frame}


\end{document}
